%\documentclass[gray,handout, pdftex, 11pt]{beamer}
%\documentclass[handout, pdftex, 11pt]{beamer}
\documentclass[pdftex, 11pt]{beamer}

\usepackage{pgfpages}
\usepackage[utf8]{inputenc}
\usepackage[T1]{fontenc}
\usepackage{lmodern}
\usepackage[italian]{babel}
\usepackage{graphicx}
\usepackage{microtype}
\usepackage{acronym}
\usepackage{array}
%\usepackage{natbib}

\usepackage{tikz}
\usetikzlibrary{intersections, arrows, shapes, decorations.pathreplacing, decorations.pathmorphing, calc}

\usepackage{appendixnumberbeamer}

\mode<presentation>{
  %-------------------------1
  \usetheme{Boadilla}
  \usecolortheme{beaver}
  %-------------------------1
  %-------------------------2
  %\usetheme{Goettingen}
  %\usecolortheme{sidebartab}
  %-------------------------2
  %\useoutertheme[right]{sidebar}
  %\usefonttheme{default}
  \setbeamercovered{transparent}
  %\setbeameroption{show notes on second screen=right}
  \setbeamertemplate{navigation symbols}{}

  \bibliographystyle{abbrv}  
  %\renewcommand\bibfont{\scriptsize}
  \setbeamertemplate{bibliography item}{\textbullet}
  \setbeamertemplate{itemize item}{\checkmark}
  \setbeamertemplate{itemize subitem}{-}
  \setbeamertemplate{enumerate items}[default]
  \setbeamertemplate{sections/subsections in toc}[square]
}

%stili
\tikzstyle{linea}=[draw, green]
\tikzstyle{service}=[very thick, draw, blue]
\tikzstyle{modulo}=[very thick, draw, red]
\tikzstyle{vettore}=[->, thick, >=stealth', draw, green]
\tikzstyle{asse}=[->, thick, >=stealth', draw, black]
\tikzstyle{lab}=[blue]
\tikzstyle{graffa}=[draw=blue, decorate,decoration={brace,amplitude=10pt}]


\newcommand{\frecciadx}{~{\tikz[baseline] \draw[freccia] (0,0.5ex) --
  (5.5ex,0.5ex);}~}


\title[Progress report]{\textbf{PROGRESS REPORT}}
\subtitle{tkLayout developers meeting}
\institute[CERN]{
  %{\Large\textbf{CERN}}\\{European Organization for Nuclear Research}\\[0.5cm]
  %\\[0.2cm]
  European Organization for Nuclear Research\\[0.5cm]
  \includegraphics[width=2cm]{img/LogoBadge.pdf}\\
}

\author[Stefano Martina]{
  %\\[0.2cm]
  \textbf{Stefano MARTINA}\\
  {\small stefano.martina@cern.ch}
}

\date[\today]{\flushright \today}

\begin{document}

\begin{frame}[plain,noframenumbering]
  \titlepage
\end{frame}

\begin{frame}
  \begin{itemize}
  \item Calcolo del \alert{$\rho$ minimo} di un poligono
    \begin{itemize}
    \item Si \alert{elimina} la dimensione \alert{$Z$} su ogni vertice
    \item Per \alert{ogni lato} del poligono proiettato si calcola:\\
      \begin{center}
      \begin{tikzpicture}[scale=0.5]
        % assi
        \path[asse] (-2cm, 0cm) to (4.5cm, 0cm) node[below]{$X$};
        \path[asse] (0cm, -1cm) to (0cm, 5cm) node[left]{$Y$};
        
        % coordinate
        \coordinate[label=above:{$P$}] (P) at (2cm, 3cm);
        \coordinate[label=above:{$R$}] (R) at (4cm, 4cm);
        \coordinate[label=above left:{$Q$}] (Q) at (-1cm, 1.5cm);
        \coordinate[label=below left:{$O$}] (O) at (0cm, 0cm);

        % modulo
        \path[modulo] (P) to (R);
        
        % vettori
        \path[vettore] (P) to (Q);
        \path[vettore] (O) to (P);
        \path[vettore] (O) to (Q);

        % graffa
        \path[graffa] (Q) to node[lab, above=3mm]{$d_P$} (P);
      \end{tikzpicture} 
      \end{center}
      \begin{align*}
        d_P &= \overrightarrow{PO} \cdot \hat{RP}&\overrightarrow{PQ} &= d_P \cdot \hat{RP}\\
        d_R &= \overrightarrow{RO} \cdot \hat{RP}&\overrightarrow{OQ} &= \overrightarrow{PQ} - \overrightarrow{OP}
      \end{align*}
      \begin{itemize}
      \item Se \alert{$d_P > 0$} e \alert{$d_R > 0$}, si prende $\rho$
        di $P$, ossia \alert{$||\overrightarrow{OP}||$}
      \item Se \alert{$d_P < 0$} e \alert{$d_R < 0$}, si prende $\rho$
        di $R$, ossia \alert{$||\overrightarrow{OR}||$}
      \item \alert{Altrimenti} $Q$ \`e all'interno del lato e si
        prende $\rho$ di $Q$, ossia \alert{$||\overrightarrow{OQ}||$}
      \end{itemize}
    \end{itemize}
  \item Ristrutturazione \alert{Usher}?
  \end{itemize}
\end{frame}

\end{document}
